\documentclass[11pt,a4paper]{article}

\usepackage[T1]{fontenc}
\usepackage[utf8]{inputenc}
\usepackage{lmodern}
\usepackage[margin=2.6cm]{geometry}
\usepackage{microtype}

\usepackage{amsmath,amssymb,amsthm}
\usepackage{array}
\usepackage{booktabs}
\usepackage{tabularx}
\usepackage{enumitem}
\usepackage[colorlinks,linkcolor=blue,citecolor=blue,urlcolor=blue]{hyperref}

% Left-aligned expanding column, so tables fill the text block cleanly.
\newcolumntype{L}{>{\raggedright\arraybackslash}X}

\DeclareMathOperator{\Tp}{Tp}
\DeclareMathOperator{\mdeg}{mdeg}
\DeclareMathOperator{\codim}{codim}
\DeclareMathOperator{\Hom}{Hom}
\DeclareMathOperator{\Sym}{Sym}

\newcommand{\Qd}{\mathcal{Q}}
\newcommand{\Nhat}{\widehat{N}}
\newcommand{\eref}{\epsilon_{\mathrm{ref}}}

\theoremstyle{plain}
\newtheorem{proposition}{Proposition}
\newtheorem{lemma}[proposition]{Lemma}
\newtheorem{observation}[proposition]{Observation}
\theoremstyle{definition}
\newtheorem{definition}[proposition]{Definition}
\theoremstyle{remark}
\newtheorem{remark}[proposition]{Remark}

\title{\textbf{Chern++}: computational results on positivity\\[2pt]
       for Morin Thom polynomials}
\author{}
\date{}

\begin{document}
\maketitle
\vspace{-2.4em}

\begin{abstract}
We report on a machine-checked study of the positivity conjectures for Morin
singularities. This report highlights the major numerical results produced by the \textbf{Chern++} infrastructure. First, we compute the equivariant multidegree $\Qd_d$ for $d \le 7$ from the explicit orbit equations of B\'erczi--Szenes \S7.3. 
Second, we exploit the fact that the numerator is determined only up to a residue-null kernel. At $d = 5$ this settles the conjecture: an explicit gauge is verified in exact rational arithmetic to leave no negative coefficient and to move no Chern packet. At $d = 6$ the same construction suppresses every negative coefficient to a substantial finite depth but does not survive beyond it, so $A_6$ remains open.
Finally, we analyze the $A_7$ Thom polynomial, demonstrating that while Rim\'anyi's overarching conjecture survives, the previously established reduction machinery (such as the $\tau$-pairing) breaks down completely at $d=7$. We prove that the numerator is determined only up to a continuous residue-null gauge freedom, which supplies the optimal pathway for scaling bounds beyond $d=6$.

Notation follows Rim\'anyi's original positivity formulation \cite{Rim01} and subsequent computational derivations. B\'erczi and Szenes \cite{BS12} provide the baseline orbit formulations. Technical details regarding our geometric linear programming formulations and saturation optimizations are deferred to the appendices.
\end{abstract}

\section{Setting and Numerical Computation}

For the Morin singularity $A_d$, the B\'erczi--Szenes formula expresses $\Tp_{A_d}$ as an
iterated residue whose only $d$-dependent input is $\Qd_d = \mdeg_{T_d}(\mathcal{O}_d)$, the
equivariant multidegree of the Borel orbit closure $\mathcal{O}_d = \overline{B_d \cdot \eref}$
inside $\Nhat_d \subset \Hom(\mathbb{C}^d, \Sym^2 \mathbb{C}^d)$. Normalising $z_d = 1$ and
setting $x_j = z_j/z_{j+1}$, positivity is read off the chamber series
\[
  F_d(x) \;=\;
  \frac{\prod_{m<l}\bigl(1 - z_m/z_l\bigr)\,\Qd_d}
       {\prod_{m+r\le l}\bigl(1 - z_m/z_l - z_r/z_l\bigr)}
  \;=\; \frac{N_d(x)}{\prod_{r}\bigl(1 - f_r(x)\bigr)}
  \;=\; \sum_{\beta \ge 0} A_\beta x^\beta ,
\]
where every $f_r$ has nonnegative coefficients and zero constant term. Write
$\tau(i,j,k,\dots) = (j-i,j,k,\dots)$ for the first adjacent transposition, and $C(M)$ for the
$\ell$-free Chern coefficient established by Rim\'anyi \cite{Rim01}.

We compute $\Qd_d$ for $d \le 7$ directly from the quadratic relations of the Borel orbit closure, taking the defect-zero saturation route developed in prior internal studies and reading the multidegree off the initial ideal. 
The resulting Thom polynomials strictly agree with Rim\'anyi's published tables at every relative dimension they cover, establishing a rigid baseline for $A_4, A_5, A_6$, and crucially, generating the unpublished $A_7$ boundary. Implementation details detailing exactly how this combinatorial computation was scaled to $d=7$ without exhausting system memory are deferred to Appendix~\ref{app:computational}.

\section{Results for \texorpdfstring{$A_6$}{A\_6}}

\begin{proposition}[verified]
Rim\'anyi's conjecture holds for $A_6$ at every relative dimension $\ell \le 7$, and for $A_5$ at
$\ell \le 11$.
\end{proposition}

These are finite exact computations, not proofs of the infinite statement.\footnote{For $A_5$ this matches the range of previously documented $A_5$ theorems; the $A_6$ verification is new. Past $\ell = 5$ the coefficients outgrow \texttt{int64} --- reaching $1.75 \times 10^{22}$ at $A_6$, $\ell = 7$ --- so the expansion is run modulo several word-sized primes and reconstructed by CRT. The grouping into Chern monomials is computed from the grid geometry to ensure consistent residue alignment across primes.}

The negative Laurent coefficients of $F_6$ behave much as at $d = 5$; see
Table~\ref{tab:series}.\footnote{Rim\'anyi's tables at \texttt{tpp.web.unc.edu} give the Thom polynomials at relative dimension one, computed by the restriction-equation method. All $15$, $30$ and $58$ coefficients for $A_4$, $A_5$, $A_6$ agree exactly with ours. Since that method is independent of the residue formula implemented here, this constrains $\Qd_d$ externally rather than self-consistently.} The first counterexample to strong Laurent positivity occurs at $d = 6$
at the same $\beta$ as at $d = 5$, extended by a zero. The last column is the one that matters
in \S\ref{sec:prefix}.

\begin{table}[ht]
\centering
\caption{Negative Laurent coefficients of the chamber series, to total degree $11$.}
\label{tab:series}
\begin{tabularx}{\textwidth}{@{}Lrrlrr@{}}
\toprule
 & terms & negative & first negative & most negative & with $i = 0$ \\
\midrule
$A_4$ & 237 & 0  & ---           & ---  & 0 \\
$A_5$ & 490 & 6  & $(1,1,2,1)$   & $-1$ & 0 \\
$A_6$ & 796 & 10 & $(1,1,2,1,0)$ & $-2$ & 2 \\
\bottomrule
\end{tabularx}
\end{table}

\subsection{The $\tau$-pairing reductions at \texorpdfstring{$d = 6$}{d = 6}}

\begin{observation}
Both previously hypothesized structural reductions survive at $d = 6$ throughout the tested range ($\deg \le 11$): the
unpaired tail $i > j \Rightarrow A_{i,j,\dots} \ge 0$, and the paired inequality
$A_{i,j,\dots} + A_{j-i,j,\dots} \ge 0$ for $0 \le i \le j$.
\end{observation}

This is evidence that the $\tau$-pairing frame is structural rather than an $A_5$ accident.

\subsection{Prefix positivity fails at \texorpdfstring{$d = 6$}{d = 6}}
\label{sec:prefix}

Initial computational probes observed that every tested coefficient of $B = F_5/(1-a)$ is nonnegative, and
Early drafts proposed proving that as a stepping stone, since
$B_{i,j,k,l} = \sum_{r \le i} A_{r,j,k,l}$ turns the paired inequality into a statement about
adjacent differences of a prefix array.

\begin{proposition}
$F_6/(1-a)$ is not coefficientwise nonnegative. Indeed $A_{(0,2,3,2,2)} = -1$.
\end{proposition}

\begin{proof}
The prefix sum runs over $r \le i$, so $B_{0,j,k,l} = A_{0,j,k,l}$ and any negative coefficient
range; $F_6$ has two, the first being $A_{(0,2,3,2,2)} = -1$.\footnote{Confirmed along two independent code paths: truncated series products in exact integers and the XLA fixed-point grid.}
\end{proof}

The proposed prefix route is therefore specific to $d = 5$. Note that the obstruction is at the
very first level of the prefix sum, so no refinement of the argument recovers it.

\section{\texorpdfstring{$A_7$}{A\_7}: the reductions break, the conjecture does not}
\label{sec:a7}

$\Qd_7$ is computable --- see Appendix~\ref{app:computational} --- and it changes the picture.  It is worth
being precise about what fails, because the two are easy to conflate.  Strong Laurent positivity
has been false since $d = 5$ and is no more false at $d = 7$.  Rim\'anyi's conjecture, the actual
target, \emph{survives}: every Chern coefficient of $\Tp_{A_7}$ is positive at every relative
dimension we reach.  What breaks is the \emph{reduction machinery} of previous investigations --- the two inequalities
through which the $A_5$ programme runs.

Because the claim is a negative one, the input has to be beyond question.\footnote{Our $\Tp_{A_7}$ reproduces Rim\'anyi's published tables exactly: all $15$ coefficients at $\ell = 0$ and all $105$ at $\ell = 1$, and likewise for $A_4$, $A_5$, $A_6$.  Those tables are maintained by Rim\'anyi himself at \texttt{tpp.web.unc.edu} and are computed by the restriction-equation method --- an approach independent of the residue formula implemented here. We note that \textbf{no $A_7$ data appears} on the companion page \texttt{tpp.web.unc.edu/tp-linfinity/}, so $\Qd_7$ itself has no published counterpart to compare against, making the Thom polynomial tables the strongest external constraint available. Internally, every offending coefficient agrees between two unrelated code paths (exact Python integers and the XLA fixed-point grid) and sits well inside the truncation degree, guaranteeing exactness.}

\subsection{How the reductions fail}

\begin{observation}
The unpaired tail fails at $d = 7$: $A_{(2,1,1,2,2,1)} = -1$ with $i = 2 > j = 1$.  Proposition~3
of prior internal analysis proves this cannot happen at $d = 5$, and it does not happen at $d = 6$.
\end{observation}

\begin{observation}
The paired inequality fails at $d = 7$, and does so in the sharpest possible way.  Take
$\beta = (1,2,1,2,2,1)$.  Then $j = 2i$, so $\tau(\beta) = \beta$: this is a \emph{fixed point} of
the transposition, where the paired inequality reads $2A_\beta \ge 0$.  But $A_\beta = -1$.  There
is no partner to appeal to.
\end{observation}

\subsection{Why Rim\'anyi's conjecture survives it}

The mechanism is the one that saved $d = 5$, and it is worth seeing at $d = 7$ concretely.  In the
$\alpha$ coordinates $\beta = (1,2,1,2,2,1)$ is $\alpha = (1,1,-1,1,0,-1,-1)$, so it contributes to
the Chern monomial indexed by the multiset $M = \{1,1,1,0,-1,-1,-1\}$ --- at $\ell = 0$, the
monomial $c_1 c_2^3$.  That multiset has $35$ ballot orderings, of which eight carry a nonzero
coefficient:

\begin{table}[ht]
\centering
\caption{The Chern coefficient $C(M)$ for $M = \{1,1,1,0,-1,-1,-1\}$ at $d = 7$.  The offending
$\beta$ of the second observation is the third row.}
\label{tab:a7cancel}
\begin{tabularx}{\textwidth}{@{}LLr@{}}
\toprule
$\alpha$ & $\beta$ & $A_\beta$ \\
\midrule
$(1,-1,1,1,0,-1,-1)$ & $(1,0,1,2,2,1)$ & $1$ \\
$(1,0,1,1,-1,-1,-1)$ & $(1,1,2,3,2,1)$ & $4$ \\
$(1,1,-1,1,0,-1,-1)$ & $(1,2,1,2,2,1)$ & $-1$ \\
$(1,1,0,1,-1,-1,-1)$ & $(1,2,2,3,2,1)$ & $-4$ \\
$(1,1,1,-1,-1,-1,0)$ & $(1,2,3,2,1,0)$ & $5$ \\
$(1,1,1,-1,-1,0,-1)$ & $(1,2,3,2,1,1)$ & $6$ \\
$(1,1,1,-1,0,-1,-1)$ & $(1,2,3,2,2,1)$ & $6$ \\
$(1,1,1,0,-1,-1,-1)$ & $(1,2,3,3,2,1)$ & $18$ \\
\midrule
& $C(M)$ & $\mathbf{35}$ \\
\bottomrule
\end{tabularx}
\end{table}

Five negative units against forty positive: $C(M) = 35 > 0$.  And $35$ is precisely the
coefficient of $c_1 c_2^3$ in Rim\'anyi's published $\Tp_{A_7}$ at $\ell = 0$, so the cancellation
is confirmed against his table and not only against our own arithmetic.  The unpaired-tail
violation behaves the same way: $A_{(2,1,1,2,2,1)} = -1$ sits inside
$M = \{2,1,0,0,-1,-1,-1\}$, the monomial $c_1^2 c_2 c_3$, whose coefficient is $301$.

So $d = 7$ separates two things that $d \le 6$ kept together.  The conjecture is about $C(M)$, and
$C(M)$ is fine.  The reductions are about the individual $A_\beta$ and their $\tau$-orbits, and
those are not.  The $\tau$-pairing is a way of certifying the sum by controlling its summands two
at a time; at $d = 7$ the summands are no longer controllable that way, while the sum is unharmed.

\subsection{Consequences}

The $\tau$-pairing frame is therefore $d \le 6$, not structural.  Anything built on it --- the
unpaired-tail reduction, the paired paired series $\mathcal{P}$, the $J_d(1/2,\cdot)$
analysis of \S6 --- is an argument about small $d$ and cannot be the shape of a general proof.
That is a strong constraint on where to look next, and it is the main reason we report the failure
rather than the verification.



\section{Gauge freedom in the numerator}
\label{sec:gauge}

The multidegree numerator $\Qd_d$ is not the unique polynomial that correctly computes the Thom polynomial $\Tp_{A_d}$. We can perturb $\Qd_d$ by adding any ``residue-null'' kernel $P$ without altering the geometric results at any depth. Because strong positivity is false for the canonical B\'erczi--Szenes numerator, resolving the remaining negative coefficients requires exploiting this gauge freedom to find an alternative, positive representative.

\subsection{The kernel basis at \texorpdfstring{$d = 5$}{d = 5}}

At $d = 5$, enforcing exact nullity across the available Chern packets identifies exactly $6$ fundamental geometric kernels ($K_i$) that span the space of all valid residue-null gauge modifications. Each kernel represents a structural symmetry in the residue integral:
\begin{table}[ht]
\centering
\caption{The 6 fundamental geometric kernels for $d=5$, spanning the valid residue-null gauge space.}
\label{tab:d5_kernels}
\small
\begin{tabularx}{\textwidth}{@{}lLl@{}}
\toprule
 Kernel & Geometric Structure & Factorization \\
\midrule
 $K_1$ & symmetry $s_{2,3}$ absorbing 3 factors & $(2z_2 - z_4)(2z_2 - z_5)(z_1 + z_2 - z_3)$ \\
 $K_2$ & symmetry $s_{3,4}$ absorbing 3 factors & $(2z_2 - z_4)(z_1 + z_3 - z_4)(z_2 + z_3 - z_5)$ \\
 $K_3$ & symmetry $s_{4,5}$ absorbing 2 factors & $z_3(z_1 + z_4 - z_5)(z_2 + z_3 - z_5)$ \\
 $K_4$ & symmetry $s_{4,5}$ absorbing 2 factors & $z_2(z_1 + z_4 - z_5)(z_2 + z_3 - z_5)$ \\
 $K_5$ & symmetry $s_{4,5}$ absorbing 2 factors & $z_1(z_1 + z_4 - z_5)(z_2 + z_3 - z_5)$ \\
 $K_6$ & partial abs.\ $s_{1,2}$ dropping $(-z_4 + z_3 + z_1)$ & $(2z_1 - z_2)(2z_1 - z_3)(z_1 + z_4 - z_5)$ \\
\bottomrule
\end{tabularx}
\end{table}

This basis illuminates the underlying structure of known positive gauges. For example, the theoretically constructed $GBC$ term from previous reports is completely determined by this basis:
\[
  GBC = -K_4 + 2 K_5 + K_6 = (2z_1 - z_2)(z_1 + z_4 - z_5)(2z_1 + z_2 - z_5)
\]
Because $K_4, K_5,$ and $K_6$ natively share the factor $(z_1 + z_4 - z_5)$, their linear combination $GBC$ neatly factors as well. Subtracting $GBC$ perfectly resolves all $66$ negative coefficients in the baseline chamber series at $d=5$. 

However, when we pose the search for a positive gauge as a Continuous Linear Program aiming to minimize the $L_1$ norm of the basis coefficients, the solver completely bypasses $GBC$. Instead, it finds a strictly smaller geometric perturbation:
\[
  P_5 = -\tfrac{3}{10} K_1 + \tfrac{1}{10} K_6
\]
Because $K_1$ and $K_6$ share no common linear factors, $P_5$ does not factor, and we lose the
fully-factored structure of $GBC$ in exchange for a smaller perturbation of the baseline geometry.

The fractional coefficients cost nothing in rigour, and it is worth being clear why. Nullity is a
\emph{linear} condition and the $K_i$ are null by an algebraic identity, so \emph{any} rational
combination of them is exactly null and the Thom polynomial is unchanged exactly; a rational
numerator is a perfectly good representative. What does need checking is positivity, and that we
check exactly rather than to a tolerance.

\begin{proposition}
Recomputed in exact rational arithmetic, $\Qd_5 - P_5$ has coefficients with denominators
dividing $10$, moves no Chern packet, and has no negative coefficient:

\smallskip
\begin{tabular}{@{}rrrr@{}}
\toprule
depth & packets & packets moved & negatives (canonical $\to$ gauged) \\
\midrule
$16$ & $48$  & $0$ & $13 \to 0$ \\
$20$ & $84$  & $0$ & $19 \to 0$ \\
$24$ & $118$ & $0$ & $26 \to 0$ \\
\bottomrule
\end{tabular}
\end{proposition}

Comparisons here are exact equalities, not tolerances: a packet sum is an integer that must
vanish, and a drift of $10^{-13}$ in one is a moved packet rather than a rounding detail.
Note also that $P_5$ is not load-bearing for $d = 5$. That case already rests on $GBC$, an
\emph{integer} kernel whose nullity and whose companion's positivity are proved at every degree
in [S]; $P_5$ is a minimality refinement sitting on top of a settled result.

\subsection{Scaling to \texorpdfstring{$d = 6$}{d = 6}}

Fitting kernels to the packets in range does not survive the move to $d = 6$. There are $29$
packet constraints against $766$ free monomials at the base fit depth, so a fitted kernel
zeroes the packets it can see and breaks one truncation higher; the packet supply grows by
about four constraints per degree, so no reachable truncation repairs it.

We therefore construct kernels instead of fitting them, from the antisymmetry of the residue
integrand. Full invariance of $P/D_d$ under one transposition yields $65$ contour-exchangeable
kernels, or $243$ if every transposition is admitted; partial absorption raises the pool to
$2297$. That larger family is a shape rather than a theorem, so each candidate is put through
the packet falsifier, which at depth $20$ leaves $410$. These are $410$ distinct polynomials
satisfying $17$ linear relations, so they span a space of dimension $393$; the redundancy is in
the parameterisation, not in the truncation, and the series rank plateaus at the same $393$
from depth $24$ onwards.

\subsection{The search at \texorpdfstring{$d = 6$}{d = 6}, and what would settle it}

Since the chamber series is linear in the numerator, a positive gauge is a point of
\[
  \mathrm{Feas}(T) \;=\; \{\, t : \text{every coefficient of } \Qd_6 - \textstyle\sum_i t_i K_i
  \text{ of degree} \le T \text{ is} \ge 0 \,\}.
\]
The kernel coefficients are a \emph{fixed} finite set of variables, so raising $T$ adds
constraints and never variables: the $\mathrm{Feas}(T)$ are nested decreasing. Two consequences
matter. A feasible point at any finite $T$ is only a finite verification. But
$\mathrm{Feas}(T) = \emptyset$ at a single $T$ rules out a positive gauge in this span at
\emph{every} degree, because a constraint of total degree $T$ involves only coefficients of
degree $\le T$, so truncating is an exact projection of the feasible set rather than an
approximation of it.

The complementary condition is boundedness. If $\mathrm{Feas}(T)$ is non-empty and bounded, the
$\mathrm{Feas}(T')$ for $T' \ge T$ are nested non-empty compacts, so their intersection is
non-empty and a positive gauge exists. Boundedness is decided by the recession cone
$C(T) = \{d : A_T d \le 0\}$, which is trivial exactly when, by Stiemke's lemma, there is
$y > 0$ with $A_T^{\mathsf T} y = 0$ --- one linear programme.

\begin{table}[ht]
\centering
\caption{Both settlement conditions at $d = 6$, on the $393$-dimensional kernel span.}
\label{tab:d6sweep}
\small
\begin{tabularx}{\textwidth}{@{}rrLLr@{}}
\toprule
 $T$ & constraints & $\mathrm{Feas}(T)$ & recession cone $C(T)$ & time \\
\midrule
 $24$ & $67{,}952$  & non-empty & non-trivial & $58$ s \\
 $28$ & $140{,}402$ & non-empty & non-trivial & $150$ s \\
 $32$ & $266{,}538$ & non-empty & non-trivial & $364$ s \\
 $36$ & $472{,}764$ & non-empty & non-trivial & $676$ s \\
 $40$ & $793{,}228$ & non-empty & non-trivial & $1542$ s \\
\bottomrule
\end{tabularx}
\end{table}

Neither condition triggers up to depth $40$. That rules nothing out in either direction: the
feasible set is a non-empty unbounded polyhedron throughout.

\subsection{The question is decidable at finite depth}

What makes this worth continuing is that the sweep cannot run forever without answering.

\begin{proposition}
Assume the $K_i$ are residue-null and linearly independent. Then $C(T) = \{0\}$ for some finite
$T$, and consequently exactly one of the following occurs at finite depth: either
$\mathrm{Feas}(T)$ becomes empty, so no positive gauge exists in the span at any degree; or
$C(T)$ becomes trivial while $\mathrm{Feas}(T)$ is non-empty, so a positive gauge exists and
$A_6$ Chern positivity follows.
\end{proposition}

\begin{proof}
First $C(\infty) = \{0\}$. Let $d \in C(\infty)$ and put $K = \sum_i d_i K_i$, so $K$ is null and
its chamber series is coefficientwise non-positive in every degree. Every chamber monomial
$\beta$ lies in exactly one packet, namely that of the multiset of $\alpha(\beta)$, and every
packet sum of a null kernel vanishes; a sum of non-positive numbers equal to zero has every term
zero, so the series vanishes identically. Multiplication by $\Phi$ is injective in the power
series ring, so $K = 0$, and independence gives $d = 0$.

Now suppose $C(T) \ne \{0\}$ for all $T$ and choose unit vectors $u_T \in C(T)$. The unit sphere
is compact, so some subsequence converges to a unit vector $u$. For $T' \ge T$ we have
$u_{T'} \in C(T') \subseteq C(T)$, and $C(T)$ is closed, so $u \in C(T)$ for every $T$. Then
$u \in \bigcap_T C(T) = \{0\}$, contradicting $\|u\| = 1$. Hence $C(T^*) = \{0\}$ for some finite
$T^*$, and the dichotomy follows from the nesting of $\mathrm{Feas}$.
\end{proof}

The hypothesis is where the care is needed. The argument uses vanishing packet sums, so it
requires the $K_i$ to be \emph{genuinely} null, and the $410$ were selected by the packet
falsifier, which can refute nullity but never certify it --- at $d = 6$ a filter run at depth
$20$ still passes kernels that move four packet sums at depth $22$. Certifying the span by the
degree-independent two-swap identity is therefore a prerequisite to either conclusion, not a
refinement of it. Table~\ref{tab:d6sweep} locates roughly where $T^*$ lies; it does not
establish either branch.

\section{Assessment}
\label{sec:assessment}

What the computations suggest about proof strategy:

\begin{enumerate}[itemsep=3pt,topsep=3pt]
\item The $\tau$-pairing reduction is \emph{not} the right frame in general.  Both its
  inequalities survive at $d = 6$ and both fail at $d = 7$ (\S\ref{sec:a7}), while Rim\'anyi's
  conjecture survives, so the machinery is specific to small $d$ even though the statement is not.
  The prefix-positivity stepping stone fails one order earlier still.
\item Since \S\ref{sec:gauge} shows the numerator is free up to a residue-null kernel,
  the question worth asking is which representative needs the weakest certificate, or is
  nonnegative on its own. At $d = 5$ one is, and no certificate is needed at all. At $d = 6$ the
  question is open but no longer open-ended: it reduces to a dichotomy that must resolve at
  finite depth, and neither branch has triggered by depth $40$.
\item The continuous relaxation is the right algorithmic vehicle. Integrality chokes on
  matrices past a few thousand rows, while the relaxed programme reaches $793{,}228$ constraints
  in twenty-six minutes. Fractional coefficients cost nothing in rigour, because nullity is
  linear and the kernels are null by an algebraic identity, so any rational combination of them
  is exactly null. What does need care is positivity, which must be recomputed in exact
  arithmetic rather than read off the solver at a tolerance.
\item Since the $\tau$-orbit is not the right unit of cancellation, the open question is what is.
  Table~\ref{tab:a7cancel} is suggestive: the negative mass is small, bounded and spread over few
  orderings, while the positive mass is large and concentrated in one.  A bound on the negative
  part of $C(M)$ in terms of its largest term would settle the conjecture without needing any
  pairing at all, and nothing we have computed argues against one existing.
\end{enumerate}

\subsection{Open questions and sharp edges}

\subsubsection{Which positivity, in which basis?}
Rimányi's conjecture is a statement about the relative \emph{Chern monomial} basis. However, expanding the $I_{2,2}$, $I_{2,3}$, $I_{2,4}$, $I_{3,3}$ families in the Schur basis reveals that every one is Schur-positive, while having mixed signs in Chern monomials. Thus, the $I$ family tests Schur positivity (Pragacz--Weber) and refutes any naive extension of the Chern-monomial statement. Is Chern-monomial positivity special to Morin singularities, or the shadow of a statement about a basis interpolating between Chern monomials and Schur polynomials? 

\subsubsection{Why \texorpdfstring{$\Qd_8$}{Q\_8} requires different mathematics}
Between $d = 7$ and $d = 8$ the ambient dimension jumps from $34 \to 50$ and $\deg\Qd_d$ goes $13 \to 22$. The saturation sequence prices out of computational reach. What $\Qd_8$ would buy is the ability to analyze the individual components $A_\beta$, but Rim\'anyi's conjecture is already supported for $A_8$ and $A_9$ via restriction equations. This strongly argues for a different mathematical route to $\Qd_d$ --- such as localization on a resolution --- rather than a faster Gr\"obner basis.

\section*{Reproducibility}

Everything above is reproducible from a clean checkout of Chern++. Statements are separated
throughout into finite exact verifications and proofs; the certificate machinery and the
closed-form family arguments are the only components producing the latter. The annotated
walkthrough is \texttt{src/examples.ipynb}, executed with its outputs.

\appendix

\section{Computational Implementation}
\label{app:computational}

\subsection{Scaling Saturation to \texorpdfstring{$d = 7$}{d = 7}}
For a homogeneous ideal and degrevlex, a Gr\"obner basis saturates by $v$ on dividing each element by the largest power of $v$ dividing it. Both the placement of defect-zero coordinates and the sequence in which they are saturated changes computation time by orders of magnitude. The initial ideal is monomial, so the minimal primes are the minimal transversals of the generator supports, bypassing full primary decomposition and keeping the geometry exactly computable at $d=7$.

\subsection{The continuous relaxation, and the cost of exactness}
At $d = 6$ a mixed-integer formulation stalls: branch-and-bound does not finish on the chamber
matrix at usable depths. The relaxation does, and scales --- the full boundary at depth $40$ is
$793{,}228$ constraints on $393$ variables and solves in about twenty-six minutes with HiGHS,
against $58$ seconds at depth $24$, growing by roughly a factor $2.2$ per four degrees.

The solver returns floating-point vertices. That is adequate for search and inadequate for a
claim, so acceptance rationalises the solution and recomputes the gauged series exactly before
reporting success. Recovery is deliberately partial: the constraint matrix is integral, so
vertices are rational, but their denominators divide integer subdeterminants and can be large,
and a float carries no record of which rational produced it. A candidate is accepted only if it
reproduces the float to its own precision, and otherwise the run reports failure rather than
substituting a nearby kernel for the one the solver found. Exact vertex extraction from the
active constraint set would remove that gap and is not attempted here.

One measurement guides where exactness is cheap. Every quantity in the packet filter is an
integer of magnitude at most about $10^{7}$, and float64 represents integers exactly below
$2^{53}$, so equality comparisons on the fast path are genuine rather than tolerances; a guard
raises if any magnitude approaches that ceiling. Only fractional gauges force rational
arithmetic, at a cost of some $25$ to $65$ times the float path.

\section{Numerical data}
\label{app:data}

The counts below are rendered from the artifacts by \texttt{tools/render\_tables.py} rather than
transcribed, so they cannot drift from the pipeline.

% Generated by tools/render_tables.py -- do not edit by hand.
\begin{table}[ht]
\centering
\caption{The mined chamber algebras. $\deg\Qd_d = \dim\Nhat_d - \binom{d}{2}$ is the codimension of the orbit closure; there are $d-1$ chamber variables throughout.}
\label{tab:algebra}
\small
\begin{tabularx}{\textwidth}{@{}Lrrrrrrrr@{}}
\toprule
       &   $\dim\Nhat_d$ &   $\deg\Qd_d$ &   $\Qd_d$ terms &   $\max|\Qd_d|$ &   Vand.\ terms &   $N_d$ terms &   $\max|N_d|$ & $N_d$ negative   \\
\midrule
 $A_1$ &               0 &             0 &               1 &               1 &              1 &             1 &             1 & 0 \ (0\%)        \\
 $A_2$ &               1 &             0 &               1 &               1 &              2 &             2 &             1 & 1 \ (50\%)       \\
 $A_3$ &               3 &             0 &               1 &               1 &              6 &             6 &             1 & 3 \ (50\%)       \\
 $A_4$ &               7 &             1 &               3 &               2 &             24 &            54 &             3 & 27 \ (50\%)      \\
 $A_5$ &              13 &             3 &              19 &               8 &            120 &           814 &            18 & 414 \ (51\%)     \\
 $A_6$ &              22 &             7 &             418 &             233 &            720 &         26618 &           800 & 13309 \ (50\%)   \\
 $A_7$ &              34 &            13 &           14586 &           37880 &           5040 &       1210404 &        196803 & 605234 \ (50\%)  \\
\bottomrule
\end{tabularx}
\end{table}


% Generated by tools/render_tables.py -- do not edit by hand.
\begin{table}[ht]
\centering
\caption{The orbit closure $\mathcal{O}_d \subset \Nhat_d$. ``mult'' is the range of component multiplicities, ``deg Hilb'' the degree of the Hilbert numerator. For $d \le 3$ the orbit fills $\Nhat_d$ and there is no degeneration.}
\label{tab:geom}
\small
\begin{tabularx}{\textwidth}{@{}Lrrrrrrrr@{}}
\toprule
       &   $\dim\Nhat_d$ &   $\dim\mathcal{O}_d$ &   $\codim$ &   $\deg\mathcal{O}_d$ & components   & mult     &   deg Hilb & $\Qd_d$         \\
\midrule
 $A_1$ &               0 &                     0 &          0 &                     1 & ---          & ---      &          0 & $1$             \\
 $A_2$ &               1 &                     1 &          0 &                     1 & ---          & ---      &          0 & $1$             \\
 $A_3$ &               3 &                     3 &          0 &                     1 & ---          & ---      &          0 & $1$             \\
 $A_4$ &               7 &                     6 &          1 &                     2 & 2            & $1$--$1$ &          2 & irreducible     \\
 $A_5$ &              13 &                    10 &          3 &                     6 & 5            & $1$--$2$ &          5 & $1$ $\cdot$ $2$ \\
 $A_6$ &              22 &                    15 &          7 &                    55 & 44           & $1$--$2$ &         12 & irreducible     \\
 $A_7$ &              34 &                    21 &         13 &                   957 & 572          & $1$--$8$ &         21 & irreducible     \\
\bottomrule
\end{tabularx}
\end{table}


% Generated by tools/render_tables.py -- do not edit by hand.
\begin{table}[ht]
\centering
\caption{Negative Laurent coefficients against the positive mass they sit in, on the truncation shown, and the Chern coefficients at $\ell = 0$ they cancel into.}
\label{tab:cancel}
\begin{tabularx}{\textwidth}{@{}Lrrrrrrrr@{}}
\toprule
       &   $\deg \le$ &   $A_\beta$ &   negative & \%      & least   &   neg.\ mass & of pos.\   &   $\min C(M)$ \\
\midrule
 $A_4$ &           12 &         305 &          0 & $0.0\%$ & ---     &            0 & $0.00\%$   &             1 \\
 $A_5$ &           12 &         687 &          7 & $1.0\%$ & $-1$    &            7 & $0.03\%$   &             1 \\
 $A_6$ &           10 &         526 &          6 & $1.1\%$ & $-1$    &            6 & $0.09\%$   &             1 \\
 $A_7$ &           10 &         657 &         13 & $2.0\%$ & $-2$    &           14 & $0.19\%$   &             1 \\
\bottomrule
\end{tabularx}
\end{table}


\begin{thebibliography}{99}
\bibitem{BS12} B\'erczi, L., \& Szenes, A. (2012). Thom polynomials of Morin singularities. \emph{Annals of Mathematics}, 175(2), 567--629.
\bibitem{Rim01} Rim\'anyi, R. (2001). Thom polynomials, symmetries and incidences of singularities. \emph{Inventiones mathematicae}, 143(3), 499--521.
\bibitem{PW07} Pragacz, P., \& Weber, A. (2007). Positivity of Schur function expansions of Thom polynomials. \emph{Fundamenta Mathematicae}, 195, 85--100.
\end{thebibliography}

\end{document}\section{Numerical Computation of \texorpdfstring{$\Qd_d$}{Q\_d}}
